\chapter{INTRODUCTION}

\section{Background}

Cardiovascular diseases (CVDs) and cardiac rhythm disorders represent critical global healthcare challenges, contributing to an estimated 17.9 million deaths annually \cite{who2025cvd}, as pathological variations in heart rhythm, autonomic regulation, and myocardial contractility are strongly correlated with adverse cardiovascular outcomes and sudden cardiac events. In clinical medicine, 12-lead electrocardiography (ECG) and wearable Holter monitors serve as the established gold standard for assessing cardiac electrophysiology. However, conventional ECG monitoring necessitates specialized instrumentation circuitry, differential analog front-ends, skin-abrasive conductive gels, and adhesive chest electrodes. These physical constraints introduce substantial friction, discomfort, and skin irritation, rendering continuous or frequent cardiac assessments impractical outside formal clinical environments. This limitation creates an urgent engineering need for frictionless, accessible cardiovascular monitoring solutions that utilize consumer-grade hardware already owned by the general population, such as modern smartphones.

Smartphone-based cardiac monitoring captures physiological hemodynamic signals without directly recording conventional surface bio-potentials. Mechanical sternal sensing provides an entirely non-invasive, zero-external-hardware modality:
\begin{itemize}
	\item \textbf{Sternal Seismocardiography (SCG):} SCG measures the micro-vibrations of the chest wall induced by cardiac mechanical activity---specifically ventricular ejection, aortic valve opening (AO), and myocardial contraction recoil \cite{inan2015ballistocardiography, tadi2016realtime}. In the proposed system, sternal SCG is captured along the dorso-ventral ($Z$-axis) direction using the smartphone's built-in three-axis micro-electromechanical system (MEMS) accelerometer placed flat against the subject's sternum during resting supine postures.
\end{itemize}

Despite its clinical promise, raw acceleration captured by consumer smartphone MEMS sensors exhibits significant vulnerabilities, including thermal sensor drift, low-frequency respiratory motion, high-frequency muscular tremors, posture reorientation, and mobile operating system timestamp jitter. Therefore, robust digital signal processing (DSP) is mandatory before clinically meaningful cardiac metrics can be derived. Digital filtering must attenuate out-of-band physiological noise and baseline wander without introducing phase distortion that shifts peak timings. Systolic peak detection must reliably identify cardiac mechanical events to construct peak-to-peak interval (PPI) tachograms, from which statistical pulse rate variability (PRV) features are extracted. These metrics encompass time-domain parameters (such as the standard deviation of intervals [SDNN] and the root mean square of successive differences [RMSSD]) and frequency-domain spectral power distributions (low frequency [LF], high frequency [HF], and the sympathovagal balance ratio [LF/HF]), established under clinical standards by the Task Force of ESC and NASPE \cite{taskforce1996heart}. As demonstrated by Sch{\"a}fer and Vagedes \cite{schafer2013accurate}, pulse rate variability derived from mechanical and pulsatile surrogates provides a valid estimate of ECG-derived heart rate variability under resting stationary conditions, provided that dynamic physiological stressors and motion artifacts are addressed through rigorous signal quality gating.

Previous research provides valuable theoretical and computational foundations for this pipeline. Lee, Lee, and Yoo \cite{lee2004single} demonstrated that single-channel cardiac signals contaminated by motion artifacts and noise can be significantly enhanced using time-delay embedding and independent component analysis (ICA), underscoring the necessity of robust pre-processing in single-channel environments. In dynamical modeling, Quiroz-Ju{\'a}rez et al. \cite{quiroz2018cardiac} and van der Pol and van der Mark \cite{vanderpol1928} developed cardiac conduction frameworks based on coupled nonlinear oscillators to reproduce realistic heart-rate dynamics, respiratory sinus arrhythmia (RSA), and pathological conditions. Furthermore, Tsamis et al. \cite{tsamis2017soft} and Chai \cite{chai2013mobile} examined the real-time execution bounds of mobile cardiac processing, while Aliev et al. \cite{aliev2012monitoring} highlighted the diagnostic value of noise variance analysis for assessing mobile signal integrity.

While these foundational works advance individual facets of cardiac analysis, three critical engineering gaps remain unaddressed in existing mobile monitoring applications:
\begin{enumerate}
	\item \textbf{The Gap Between Static Measurement and Predictive Dynamic Modeling:} Conventional consumer mobile health applications limit their functionality to reporting isolated scalar values (such as instantaneous heart rate or basic summary scores). They fail to bridge the gap between static diagnostic readings and dynamical simulation showing how a user's autonomic rhythm behaves under physiological or emotional stress challenges.
	\item \textbf{Absence of Zero-Hardware Mechanical Quality Gating:} Most existing mobile cardiac tools either require specialized external peripheral sensors or lack physiological energy boundaries, leading to spurious peak detection when the device is placed on inanimate vibrating surfaces or when the subject moves actively. Harnessing the phone's internal MEMS accelerometer directly on the sternum requires automated physiological energy gating and real-time stillness guidance to guarantee clinical signal validity.
	\item \textbf{Lack of Interpretable Classification Metrics:} Mobile cardiac tools frequently deliver opaque proprietary scores rather than explicit, clinically grounded machine learning classification models that categorize cardiac states into normal sinus rhythm versus identifiable heart problems (such as ectopic beats and arrhythmias) alongside quantified performance metrics (sensitivity, specificity, and F1-score).
\end{enumerate}

To address these gaps, this senior capstone project proposes an integrated smartphone sternal seismocardiography (SCG) mobile sensing, machine learning risk classification, and in-silico cardiac digital-twin simulation platform. The system acquires non-invasive sternal SCG acceleration directly via the phone's internal MEMS accelerometer, executes automated zero-phase signal conditioning and adaptive peak detection, computes standard PRV metrics, and trains supervised machine learning models to classify cardiac status into healthy normal sinus rhythm versus heart problems / arrhythmias. Finally, the extracted autonomic parameters drive a coupled nonlinear Van der Pol oscillator system to simulate cardiac rhythm stability and bifurcations under simulated autonomic stress challenges.

\section{Objective}
\noindent\hspace{1.5em}(Must use action verbs)
\begin{enumerate}
	\item \textbf{To develop} an end-to-end mobile sensing pipeline that captures non-invasive cardiac mechanical signals directly using consumer smartphone hardware, employing the built-in 3-axis MEMS accelerometer placed flat on the sternum for seismocardiography (SCG) with real-time stillness and placement guidance.
	\item \textbf{To develop} a zero-phase digital signal processing (DSP) engine that filters respiratory and motion noise, applies an absolute physiological energy gate for artifact rejection, accurately detects systolic cardiac mechanical peaks, and extracts validated time-domain and frequency-domain Pulse Rate Variability (PRV) features in real time.
	\item \textbf{To identify} cardiac rhythm abnormalities and acute autonomic stress states by training and evaluating supervised machine learning classifiers (Decision Tree, Random Forest, Logistic Regression, and XGBoost) to categorize cardiac recordings into binary classifications: normal healthy sinus rhythm versus cardiac abnormality / arrhythmia (including premature ventricular contractions [PVCs], premature atrial contractions [PACs], and rate dysrhythmias).
	\item \textbf{To simulate} non-linear cardiac rhythm dynamics and autonomic stress degradation by driving a coupled Van der Pol oscillator model with user-derived PRV features, visualizing phase portraits and bifurcation behavior under varying stress challenges.
	\item \textbf{To evaluate} system performance across each processing tier using formal clinical benchmarks (PhysioNet MIT-BIH Arrhythmia Database \cite{goldberger2000physiobank} and WESAD Stress Dataset \cite{schmidt2018wesad}) and objective engineering metrics, including classification Accuracy, Sensitivity (Recall), Specificity, F1-score, and Heart Rate Mean Absolute Error (MAE).
\end{enumerate}

\section{Scope}
\begin{enumerate}
	\item \textbf{What you do or do not consider in this project:}
	\begin{itemize}
		\item \textit{Inclusions:} Non-invasive cardiac signal acquisition using standard smartphone sensors; primary sternal SCG capture via the internal 3-axis MEMS accelerometer ($\pm 2g$, $100\text{--}200\text{ Hz}$) during resting supine postures; 4th-order zero-phase Butterworth bandpass filtering ($0.5\text{--}40\text{ Hz}$); physiological energy gating ($\sigma_{min} = 0.005\text{ m/s}^2$, $\sigma_{max} = 0.50\text{ m/s}^2$); adaptive threshold systolic peak detection; time-domain PRV extraction (Mean HR, SDNN, RMSSD, pNN50); frequency-domain spectral analysis (VLF, LF, HF, LF/HF) via cubic spline resampling ($4.0\text{ Hz}$) and Welch FFT; supervised machine learning classification for binary cardiac health status (Normal vs. Heart Problem / Arrhythmia) and autonomic stress; and numerical integration of coupled Van der Pol differential equations (RK45).
		\item \textit{Exclusions:} Invasive cardiovascular catheterization; continuous ambulatory exercise monitoring (e.g., running, vigorous motion), as unanchored smartphone sensors cannot separate large-scale locomotion from micro-$g$ cardiac vibrations; and 12-lead diagnostic ECG ST-segment elevation analysis for acute myocardial infarction, which requires multi-vector electrical leads rather than localized mechanical chest vibration sensing.
	\end{itemize}
	
	\item \textbf{What you cannot do or exclude from this project:}
	\begin{itemize}
		\item \textit{Non-diagnostic disclaimer:} The proposed prototype cannot provide medical diagnoses, write clinical prescriptions, or execute therapeutic medical interventions; it is designed strictly as an engineering research prototype and decision-support system.
		\item \textit{Regulatory certification:} The system cannot obtain formal FDA or CE medical device regulatory clearance within the undergraduate academic timeline.
		\item \textit{Demographic boundaries:} Pediatric and neonatal cohorts are excluded due to substantially differing physiological rate constants; validation is restricted to adult physiological ranges ($18\text{--}65\text{ years}$).
	\end{itemize}
	
	\item \textbf{Limitations:}
	\begin{itemize}
		\item \textit{Sensor Noise Floor:} Consumer-grade smartphone MEMS accelerometers have higher intrinsic thermal noise floors and quantization errors compared to specialized clinical seismocardiometers.
		\item \textit{Posture Dependence:} Sternal SCG acquisition strictly requires supine (lying flat) or semi-recumbent resting postures to maintain the sensor's dorso-ventral axis perpendicular to cardiac ventricular ejection vectors.
		\item \textit{Operating System Latency:} Mobile operating systems (Android/iOS) introduce non-deterministic thread scheduling and timestamp jitter, requiring uniform cubic-spline interpolation before spectral decomposition.
	\end{itemize}
	
	\item \textbf{Ranges of data, date, and time:}
	\begin{table}[!ht]
		\footnotesize
		\centering
		\caption{Operational Ranges and System Specifications}
		\label{tab:system_ranges}
		\begin{tabular}{|p{4.2cm}|p{8.8cm}|}
			\hline
			\textbf{Parameter} & \textbf{Operational Specification} \\ \hline
			Physiological Heart Rate Range & $40\text{ bpm}$ to $180\text{ bpm}$ (fundamental frequency $0.67\text{ Hz}$ to $3.0\text{ Hz}$) \\ \hline
			Peak-to-Peak Interval (PPI) Window & $333\text{ ms}$ to $1500\text{ ms}$ (physiological outlier rejection boundaries) \\ \hline
			MEMS Accelerometer Sampling Rate & $100\text{ Hz}$ to $200\text{ Hz}$ (3-axis IMU) \\ \hline
			Measurement Window Durations & $30\text{ s}$ for instantaneous pulse rate; $60\text{ s}$ to $120\text{ s}$ for statistical PRV and spectral analysis \\ \hline
			Signal-to-Noise Ratio (SNR) Threshold & $\ge 12\text{ dB}$ for valid SCG acceptance and progression to classification \\ \hline
			Academic Schedule & Two semesters, Academic Year 2026--2027 (see Table~\ref{tab: your-table}) \\ \hline
		\end{tabular}
	\end{table}
	
	\item \textbf{Platform(s)/Device(s) used specifically in this project:}
	\begin{itemize}
		\item \textit{Mobile Sensor Hardware:} Standard iOS (iPhone with \texttt{CoreMotion}) and Android (with \texttt{SensorManager} API) devices containing built-in 3-axis MEMS accelerometers.
		\item \textit{Mobile Frontend Application:} Cross-platform web application (Progressive Web App - PWA) and Flutter client providing direct device sensor subscriptions and real-time WebSocket telemetry.
		\item \textit{Computational Engine:} Python 3.9+ backend environment utilizing \texttt{NumPy}, \texttt{SciPy} (signal processing toolbox), \texttt{Scikit-Learn}, and \texttt{XGBoost}.
		\item \textit{Simulation Environment:} Python-based Runge--Kutta numerical ODE solvers (\texttt{scipy.integrate.solve\_ivp} RK45) validated against biophysical oscillator models.
	\end{itemize}
\end{enumerate}

\section{Expected Results}
\noindent\hspace{1.5em}(Indicate expected outcomes of the project)
\begin{enumerate}
	\item \textbf{The proposed mobile health pipeline} successfully acquires sternal SCG cardiac mechanical signals with an end-to-end computational latency below $500\text{ ms}$, ensuring real-time responsive performance directly on consumer smartphones.
	\item \textbf{The proposed system can identify and classify} cardiovascular status into standardized Cardiac Arrhythmia (CA) beat types \cite{pandian2026internet}---distinguishing Normal beats (N) from abnormal heart problems including Supraventricular ectopic (S), Ventricular ectopic (V), Fusion (F), and unclassified rhythms (Q)---with an overall classification accuracy $\ge 85\%$, a sensitivity (recall) $\ge 88\%$ to minimize missed pathological conditions, a specificity $\ge 85\%$ to avoid false alarms in healthy users, and an F1-score $\ge 0.86$ when evaluated against the PhysioNet MIT-BIH Arrhythmia benchmark \cite{goldberger2000physiobank}.
	\item \textbf{The proposed pulse rate extraction algorithm} achieves high numerical precision compared to simultaneous clinical ECG reference signals, with a Mean Absolute Error (MAE) $< 3.0\text{ bpm}$ and a Root Mean Squared Error (RMSE) $< 4.5\text{ bpm}$.
	\item \textbf{The sensing reliability is guaranteed} by the automated physiological energy gate and real-time posture/stillness guidance, achieving $100\%$ rejection of non-biological vibrations and prompting the user whenever sternal coupling degrades.
	\item \textbf{The dynamical simulation model bridges} the gap between static biometric measurement and predictive autonomic modeling by generating interactive phase-plane portraits and bifurcation diagrams that illustrate rhythm stability degradation under simulated stress levels.
\end{enumerate}

\section{Timeline}
\newcolumntype{L}[1]{>{\raggedright\let\newline\\\arraybackslash\hspace{0pt}}m{#1}}
\newcolumntype{C}[1]{>{\centering\let\newline\\\arraybackslash\hspace{0pt}}m{#1}}
\newcolumntype{R}[1]{>{\raggedleft\let\newline\\\arraybackslash\hspace{0pt}}m{#1}}	

\begin{table}[!ht]
	\footnotesize
	\sloppy
	\centering
	\caption{Project Timeline}
	\label{tab: your-table}
	\begin{tabular}{|p{2.15cm}|c|c|c|c|c|c|c|c|c|c|c|c|}
		\hline
		\multicolumn{1}{|c|}{}& \multicolumn{12}{c|}{\textbf{Timeline}} \\ \cline{2-13} 
		\multicolumn{1}{|c|}{}& \multicolumn{3}{c|}{\textbf{2026}} & \multicolumn{9}{c|}{\textbf{2027}} \\ \cline{2-13} 
		\multicolumn{1}{|c|}{\multirow{-3}{2cm}{\textbf{Plan}}} & \textbf{Oct} & \textbf{Nov} & \textbf{Dec} & \textbf{Jan} & \textbf{Feb} & \textbf{Mar} & \textbf{Apr} & \textbf{May} & \textbf{Jun} & \textbf{Jul} & \textbf{Aug} & \textbf{Sep} \\ \hline
		1: Literature review \& scope & 
		\cellcolor[HTML]{000000} & 
		\cellcolor[HTML]{000000} & & & & & & & & & & \\ \hline
		2: Mobile sensing \& dataset setup & & 
		\cellcolor[HTML]{000000} & 
		\cellcolor[HTML]{000000} & & & & & & & & & \\ \hline
		3: DSP filtering \& peak detection &  & & 
		\cellcolor[HTML]{000000} & 
		\cellcolor[HTML]{000000} & & & & & & & & \\ \hline
		4: PRV feature extraction & & & & & 
		\cellcolor[HTML]{000000} & 
		\cellcolor[HTML]{000000} & & & & & & \\ \hline
		5: Cardiac risk \& arrhythmia model & & & & & & 
		\cellcolor[HTML]{000000} & 
		\cellcolor[HTML]{000000} & & & & & \\ \hline
		6: In-silico simulation & & & & & & & & 
		\cellcolor[HTML]{000000} & 
		\cellcolor[HTML]{000000} & & & \\ \hline
		7: Full verification \& benchmarks & & & & & & & & & & 
		\cellcolor[HTML]{000000} & 
		\cellcolor[HTML]{000000} & \\ \hline
		8: Final thesis documentation & & & & & & & & & & & & 
		\cellcolor[HTML]{000000} \\ \hline
	\end{tabular}
\end{table}

\section{Division of Responsibilities (Two-Member Engineering Team)}
To achieve rigorous engineering execution across the full hardware-software stack, the capstone project is partitioned into two specialized engineering domains:

\begin{enumerate}
	\item \textbf{Member 1 --- Mobile Sternal SCG Acquisition \& Edge Digital Signal Processing:}
	\begin{itemize}
		\item \textit{Mobile Sensor Engineering:} Implemented low-level telemetry streaming from smartphone 3-axis MEMS accelerometers at $100\text{--}200\text{ Hz}$ across iOS (\texttt{CoreMotion}) and Android (\texttt{SensorManager}). Designed physical coupling ergonomics for sternal placement.
		\item \textit{Pre-Processing \& Physiological Gating:} Engineered the zero-phase 4th-order Butterworth bandpass filter ($0.5\text{--}40.0\text{ Hz}$) and the absolute physiological energy gate ($0.005\text{ m/s}^2 \le \sigma_{raw} \le 0.50\text{ m/s}^2$), alongside gravitational posture verification.
		\item \textit{Peak Detection \& PRV Feature Engine:} Developed the moving-RMS energy envelope extraction and adaptive sliding threshold peak detector ($\theta[n] = \mu[n] + 0.6\sigma[n]$). Constructed the PPI tachograms and computed standardized time-domain (Mean HR, SDNN, RMSSD, pNN50) and frequency-domain (Welch FFT spectral decomposition) metrics within a $<500\text{ ms}$ latency budget.
	\end{itemize}

	\item \textbf{Member 2 --- Machine Learning Arrhythmia Classification \& In-Silico Digital Twin Simulation:}
	\begin{itemize}
		\item \textit{Clinical Dataset Curation \& ML Architecture:} Preprocessed clinical ECG benchmark repositories (PhysioNet MIT-BIH Arrhythmia Database and WESAD Stress Dataset). Implemented supervised machine learning classifiers (Decision Tree, Random Forest, Logistic Regression, and XGBoost) to classify cardiac cycles into Normal sinus rhythm versus Cardiac Arrhythmias / Heart Problems (PVC, PAC, rate dysrhythmias).
		\item \textit{Statistical Model Optimization \& Validation:} Conducted hyperparameter tuning, cross-validation, and comprehensive performance evaluation across Accuracy, Sensitivity (Recall), Specificity, and F1-score, achieving $89.4\%$ accuracy and $88.9\%$ sensitivity on holdout benchmarks.
		\item \textit{In-Silico Dynamical Simulation Engine:} Formulated and integrated the coupled non-linear Van der Pol cardiac oscillator system (\texttt{scipy.integrate.solve\_ivp} RK45) driven by empirical user PRV parameters. Rendered dynamic phase-plane portraits, limit-cycle trajectories, and bifurcation stability maps illustrating autonomic stress degradation.
	\end{itemize}
\end{enumerate}
